\documentclass{homework}
% \author{Tashfeen, Ahmad}  % uncomment with your name
\class{CSCI 3223: Tashfeen's Machine Learning II}
\title{Homework 4}
\address{
  Oklahoma City University,
  Petree College of Arts \& Sciences,
  Computer Science
}

\usepackage[plain]{algorithm}
\DeclareRobustCommand{\DO}{\mathbf{do\;}}
\DeclareRobustCommand{\IF}{\mathbf{if\;}}
\DeclareRobustCommand{\THEN}{\mathbf{then\;}}
\DeclareRobustCommand{\WHILE}{\mathbf{while\;}}
\DeclareRobustCommand{\FOR}{\mathbf{for\;}}
\DeclareRobustCommand{\SL}{\mathbf{Select}}
\DeclareRobustCommand{\FIT}{\mathbf{Fitness}}
\DeclareRobustCommand{\CROSS}{\mathbf{Cross}}
\DeclareRobustCommand{\MUTATE}{\mathbf{Mutate}}

\begin{document} \maketitle

\textit{Genetic Algorithms} (GA) mimic biological evolution and
natural selection to address mathematical optimisation problems. They
begin with an initial population $P$ of solutions and aim to optimize
them over several generational iterations, similar to how desirable
traits emerge in nature (or rather, undesirable traits go extinct).

To solve an optimisation problem with a GA, we start with a
\emph{genotype} that encodes a solution. Genotype instances are
called \emph{individuals} $x$. A set of individuals is a
\emph{population} $x \in P$. Populations of individuals evolve
through \emph{generations} via reproduction. Using a
\emph{selection} $\SL(P, \FIT)$ process based on a
\emph{fitness-function} $\FIT: P \ra [0, 1]$, two parent
individuals $a, b$ are selected for \emph{crossover} to produce a
child individual $t \leftarrow \CROSS(a, b)$ in the next
generation $P_\text{next}$. The child $t$ may $\MUTATE(t)$. See
Algorithm.~\ref{ga}.

\begin{algorithm}
  \begin{flushleft}
    \noindent\textbf{Input}:
    \begin{enumerate}
      \item Initial population: $P$
      \item Fitness-function: $\FIT: P \ra [0, 1]$
      \item Error threshold $\ep$
      \item Mutation probability
    \end{enumerate}
    \noindent\textbf{Output}:
    The best individual in population, according to the fitness.
  \end{flushleft}%
  \begin{enumerate}[label=\arabic*:,noitemsep,partopsep=0pt,topsep=0pt,parsep=0pt]%
    \item $\DO$
    \item $\quad P_\text{next} \leftarrow \nil$
    \item $\quad \FOR (a, b) \leftarrow \SL(P, \FIT)$
    \item $\quad\quad t \leftarrow \CROSS(a, b)$
    \item $\quad\quad \IF (\text{small random probability})\;\THEN t \leftarrow \MUTATE(t)$
    \item $\quad\quad P_\text{next} \leftarrow P_\text{next} \cup \{t\}$
    \item $\quad P \leftarrow P_\text{next}$
    \item $\WHILE \forall x \in P, \FIT(x) < (1 - \ep)$
  \end{enumerate}%
  \caption{%
    Pseudo code for a simple Genetic Algorithm. See Mitchell, pg. 251
    or Russell and Norvig, pg 129 for more details.
  }%
  \label{ga}%
\end{algorithm}%

\question\label{example} Watch,
\begin{itemize}
  \item \href{https://www.youtube.com/watch?v=MacVqujSXWE}{The Knapsack Problem \& Genetic Algorithms - Computerphile}
  \item \href{https://www.youtube.com/watch?v=jPcBU0Z2Hj8}{The 8 Queen Problem - Numberphile}
\end{itemize}

and answer,

\begin{enumerate}[label=(\alph*)]
  \item For $S = \{1, 5, 10, 25\}, m = 31$, give a subset of $S$ that
        sums to $m$.
  \item Give a possible solution to the 8 queens problem.
\end{enumerate}

\question\label{knpga} Design a GA for the Knapsack Problem where $\abs{S}=d$ by
stating the genotype, strategy for the initial population, fitness
function, selection process, mutation and a crossover operator.

\question In the worst possible configuration of the 8 queens,
\begin{enumerate}
  \item How many queens can attack a single queen assuming attacks are not
        blocked by intermediate queens?
  \item How many total attacks (on each queen) are possible on the board?
\end{enumerate}

\question\label{plan} Design a Genetic Algorithm to approximate
solutions to the 8 queens problem by stating the genotype, strategy
for the initial population, fitness function, selection process,
mutation and a crossover operator.

\question\label{code} Using the strategies you came up with in
question \ref{plan}, finish the code given in the listing
\href{https://tashfeen.org/s/mll/stubs.py}{available here}. Note
that you'll need to finish the functions as well as use them in your
implementation of the genetic algorithm with the following
hyper-parameters,

\begin{itemize}
  \item Population size: 100
  \item Muration rate: 0.005
  \item Stopping criteria: Stop when the difference between the current
        population's fitness and the one before is less than $0.0001$. The
        fitness of a population is the mean fitness of its individuals.
\end{itemize}

Put your code for this question in \texttt{queens.py}.

\begin{enumerate}
  \item Plot the fitness against generations from question \ref{code}.
  \item What is the fitness of the fittest individual in your final
        generation? Also give the corresponding chessboard.
\end{enumerate}

\question If you've encoded the board state in an efficient way, you'll be
able to efficiently brute force search the whole space of
individuals and calculate their fitness. The individuals $x$ with
fitness $f(x) = 1$ are exact solutions to the 8 queens problem.
How many such solutions are there?

\question Generate a random knapsack problem where,
\begin{itemize}
  \item An element $s_i \in S$ is uniformly randomly sampled between $1$ and $2^{32}$.
  \item The cardinality of the solution subset $\chi$ that sums to $m$
  is at least a third of the cardinality of $S$.
  \[
    \abs{\chi} \geq \abs{S}/3
  \]
\end{itemize}

Implement your GA from question \ref{knpga} to find the best
approximation of $\chi$ and put the code in \texttt{subset.py}.

\begin{enumerate}
  \item Plot the fitness against generations.
  \item Let the subset $\hat{\chi}$ correspond to the fittest
  individual in your final generation, give $\floor{\log_2(m)} + 1$
  and $\floor{\log_2\paren{\abs{m - (\sum_{b \in \hat{\chi}} b)}}} + 1$.
\end{enumerate}

\section*{Submission Instructions}
\begin{enumerate}
  \item Submit a PDF that answers all the questions including any plots.
  \item Submit the Python files \eg \texttt{quuens.py} and \texttt{subset.py}.
\end{enumerate}
\end{document}
